% =============================================================================
% SECTION 2: MODEL
% =============================================================================

\section{The Model}

This section presents a formal model of firm compliance decisions in a compute permit market with imperfect monitoring. We derive the conditions under which firms choose to comply with permit requirements and characterize how regulatory parameters affect aggregate compliance.

\subsection{Setup}

Consider a market with $N$ firms, indexed by $i \in \{1, \ldots, N\}$. Each firm seeks to conduct a training run requiring compute $C_i$, measured in FLOP. A regulatory threshold $\tau$ (e.g., $10^{25}$ FLOP) determines which training runs require permits. Firms with $C_i > \tau$ must hold a valid permit to operate legally.

Each firm has a private valuation $v_i$ representing the economic benefit of completing its training run. This valuation captures expected revenues from deploying the trained model, strategic advantages from technological leadership, and other sources of value. Firms are heterogeneous in their valuations, reflecting differences in business models, market positions, and technical capabilities.

Firms face a binary choice: comply with permit requirements or defect by operating without a permit. A compliant firm purchases the required permits at market price $P$ and operates legally. A defecting firm avoids permit costs but faces the risk of detection and penalty.

\subsection{The Deterrence Condition}

A firm's decision to comply depends on comparing the cost of compliance against the expected cost of defection. Let:

\begin{itemize}
    \item $g_i$ denote the gain from defection for firm $i$.
    \item $p_{\mathrm{detect},i}$ denote the effective probability of detecting firm $i$'s violation.
    \item $B_i$ denote the total perceived penalty upon detection.
\end{itemize}

A risk-neutral firm complies if and only if the expected penalty from defection exceeds the gain:

\begin{equation}
p_{\mathrm{detect},i} \times B_i \geq g_i
\end{equation}

This is the central deterrence condition of the model. When it holds, compliance is the economically rational choice. When it fails, firms have an incentive to defect.

The gain from defection for firm $i$ is:

\begin{equation}
g_i = \min(p^*, v_i) + r_i \, V_i
\end{equation}

where $p^*$ is the current market clearing price for permits, $v_i$ is the firm's private valuation, $V_i$ is the capability value of running a frontier model, and $r_i$ is a firm-specific racing factor capturing competitive urgency. The $\min$ operator reflects that a firm whose valuation falls below the permit price would not have purchased a permit regardless, so the savings from evasion are bounded by $v_i$. In most calibrations $r_i \, V_i = 0$, so the racing term drops out and the gain reduces to the lesser of the permit price and the firm's valuation.

\subsection{Detection Probability}

Detection follows a two-stage sequential process. In the first stage, firm $i$ is selected for audit with probability $p_{\mathrm{audit}}(i)$. In the second stage, conditional on an audit occurring, the violation is detected with probability $p_{\mathrm{catch}}$. The overall detection probability is therefore:

\begin{equation}
p_{\mathrm{detect}}(i) = p_{\mathrm{audit}}(i) \times p_{\mathrm{catch}}
\end{equation}

The audit is the sole mechanism that can produce a legally demonstrable finding. Whistleblower reports, hardware monitoring, and historical backchecks serve as information sources that improve the probability of detection \emph{within} an audit, rather than as independent enforcement channels. This reflects the institutional reality that evidence gathered outside a formal audit process typically cannot, on its own, establish a regulatory violation.

The audit probability depends on the regulatory mode. Under random auditing, all firms face the base audit rate:

\begin{equation}
p_{\mathrm{audit}}(i) = \pi_0
\end{equation}

Under signal-dependent targeting, the audit probability also reflects the strength of the suspicion signal generated by the firm's unpermitted excess compute:

\begin{equation}
p_{\mathrm{audit}}(i) = \min\!\Big(1,\; \pi_0 + c(i) \cdot s(i) \cdot (1 - \pi_0)\Big)
\end{equation}

where $c(i) \geq 0$ is a firm-specific audit coefficient reflecting visibility or regulatory scrutiny, and $s(i) \in [0,1]$ is the suspicion signal (defined in Section~\ref{sec:enforcement}).

Conditional on an audit, detection proceeds through four nested channels in sequence: a direct audit pass, a historical backcheck, whistleblower disclosure, and hardware monitoring. The violation is caught if any channel fires:

\begin{equation}
p_{\mathrm{catch}} = 1 - \beta \cdot (1 - p_b) \cdot (1 - p_w) \cdot (1 - p_m)
\end{equation}

where $\beta \in [0,1]$ is the false negative rate of the direct audit pass, $p_b$ is the backcheck probability, $p_w$ is the whistleblower detection probability, and $p_m$ is the monitoring detection probability. The miss probability is the joint probability that all four channels fail. This sequential nesting also allows each detection event to be attributed to a single channel in the simulation output.

\subsection{Penalty Structure}

The total perceived penalty $B_i$ upon detection is firm-specific:

\begin{equation}
B_i = (\phi + K + R_i) \times \rho_i
\end{equation}

where $\phi$ is the monetary fine imposed upon violation, $K$ is the collateral forfeited, $R_i$ is firm $i$'s perceived reputational cost if caught, and $\rho_i \geq 0$ is a risk profile multiplier scaling the firm's sensitivity to losses. Collateral is posted upfront before the firm operates and is returned if no violation is detected. This structure addresses the limited liability problem: firms with insufficient assets to pay fines can still be deterred if they must post collateral in advance. The reputation term $R_i$ introduces heterogeneity in how firms perceive the consequences of detection beyond direct financial penalties, and is the mechanism through which repeated violations compound deterrence in dynamic scenarios (see Appendix~D).

\subsection{Equilibrium Compliance}

Given the deterrence condition, aggregate compliance has two components. Firms holding enough permits to cover their full compute comply automatically: the permit purchase itself constitutes compliance. Only firms with unpermitted excess compute evaluate the deterrence condition. The compliance rate is therefore:

\begin{equation}
\text{Compliance Rate} = \frac{|\mathcal{P}|}{N} + \frac{|\mathcal{U}|}{N} \cdot \Pr\!\big(p_{\mathrm{detect},i} \times B_i \geq g_i \;\big|\; i \in \mathcal{U}\big)
\end{equation}

where $\mathcal{P}$ is the set of fully permitted firms and $\mathcal{U}$ is the set of firms with unpermitted excess compute. In supply-constrained scenarios where $Q < N$, the first term is bounded by $Q/N$ and the second term captures any additional compliance induced by enforcement. In universal-access scenarios where $Q = N$ and all firms can obtain permits, the distinction collapses and the compliance rate depends entirely on whether permit purchase is individually rational.

This expression clarifies the levers available to regulators. Compliance can be increased by:

\begin{enumerate}
    \item Increasing $p_{\mathrm{detect},i}$ through higher audit rates, better monitoring, or whistleblower programs.
    \item Increasing $B_i$ through higher fines, collateral requirements, or reputation costs.
    \item Decreasing $g_i$ through more efficient permit markets that reduce the cost of compliance.
\end{enumerate}

The smart enforcement approach emphasizes the third lever: by making permits cheaper and more accessible, regulators reduce the gain from defection, allowing compliance to be maintained even with moderate detection probabilities.

\subsection{Model Parameters}

Table \ref{tab:parameters} summarizes the key parameters of the model. Detailed calibration values for each regulatory scenario are provided in Appendix A.

\begin{table}[h]
\centering
\caption{Key Model Parameters}
\label{tab:parameters}
\begin{tabular}{@{}llp{7cm}@{}}
\toprule
\textbf{Parameter} & \textbf{Symbol} & \textbf{Description} \\
\midrule
Base audit rate & $\pi_0$ & Probability any firm is audited each period \\
Audit coefficient & $c(i)$ & Firm-specific signal scaling (signal mode only) \\
False negative rate & $\beta$ & Probability direct audit pass misses a violation \\
Backcheck probability & $p_b$ & Historical review catch rate (within audit) \\
Whistleblower probability & $p_w$ & Disclosure-based catch rate (within audit) \\
Monitoring probability & $p_m$ & Hardware/metering catch rate (within audit) \\
Penalty amount & $\phi$ & Fine imposed upon detected violation \\
Collateral & $K$ & Upfront deposit forfeited upon violation \\
Reputation sensitivity & $R_i$ & Perceived brand/trust cost if caught \\
Risk profile & $\rho_i$ & Multiplier on perceived losses \\
Capability value & $V_i$ & Strategic value of running a frontier model \\
Racing factor & $r_i$ & Urgency multiplier on capability value \\
FLOP threshold & $\tau$ & Compute level above which permits are required \\
\bottomrule
\end{tabular}
\end{table}